\documentclass[../main.tex]{subfiles}
\begin{document}

\section{其他}

\subsection{类欧几里得算法}

\subsubsection{直线下整点个数}

\begin{frame}{直线下整点个数}
求解如下函数

$$
f(a,b,c,n)=\sum_{i=0}^n\left\lfloor\frac{ai+b}{c} \right\rfloor
$$\pause

实际上这是类欧几里得算法的一部分。由于形式简单，应用相对较多，故单独讲解。
\end{frame}

\begin{frame}
若 $a\ge c$ 或 $b\ge c$，可以转化为 $a<c,b<c$ 的情况：

$$
\begin{aligned}
&f(a,b,c,n)=\sum_{i=0}^n\lfloor\frac{ai+b}{c}\rfloor\\
&=\sum_{i=0}^n\lfloor\frac{(\lfloor\frac{a}{c}\rfloor c+(a\bmod c))i+\lfloor\frac{b}{c}\rfloor c+(b\bmod c)}{c}\rfloor\\
&=\frac{n(n+1)}{2}\lfloor\frac{a}{c}\rfloor+(n+1)\lfloor\frac{b}{c}\rfloor+\sum_{i=0}^n\lfloor\frac{(a\bmod c)i+(b\bmod c)}{c}\rfloor\\
&=\frac{n(n+1)}{2}\lfloor\frac{a}{c}\rfloor+(n+1)\lfloor\frac{b}{c}\rfloor+f(a\bmod c,b\bmod c,c,n)
\end{aligned}
$$
\end{frame}

\begin{frame}
接下来考虑一般情况 $a<c$ 且 $b<c$：

$$
\begin{aligned}
f(a,b,c,n)&=\sum_{i=0}^n\lfloor\frac{ai+b}{c}\rfloor\\
&=\sum_{i=0}^n\sum_{j=0}^{\lfloor\frac{ai+b}{c}\rfloor-1}1\\
&=\sum_{j=0}^{\lfloor\frac{an+b}{c}\rfloor-1}\sum_{i=0}^n[j<\lfloor\frac{ai+b}{c}\rfloor]\\
\end{aligned}
$$\pause

$$
j<\lfloor\frac{ai+b}{c}\rfloor\iff\frac{jc+c-b-1}{a}<i
$$
\end{frame}

\begin{frame}
令 $m=\lfloor\frac{an+b}{c}\rfloor$，有

$$
\begin{aligned}
f(a,b,c,n)&=\sum_{j=0}^{m-1}n-\lfloor\frac{jc+c-b-1}{a}\rfloor\\
&=nm-f(c,c-b-1,a,m-1)
\end{aligned}
$$\pause

于是可以递归计算。复杂度 $O(\log n)$。
\end{frame}

\subsubsection{}

\begin{frame}{CF1912J Joy of Pok´emon Observation}
给定 $l_1,l_2,l_3,t$，求满足 $l_1x+l_2y+l_3z=t$ 的非负整数对 $(x,y,z)$ 的数量。

$l_1,l_2,l_3\le 16,t\le 10^9$。
\end{frame}

\begin{frame}
将 $x,y$ 在 $\bmod l_3$ 意义下考虑，设 $x=pl_3+x',y=ql_3+y'$，则 $l_3(l_1p+l_2q+z)=t-l_1x'-l_2y'$。\pause

枚举 $x',y'$，若 $l_3\nmid  t-l_1x'-l_2y'$ 则无解，否则记 $K=\dfrac{t-l_1x'-l_2y'}{l_3}$，变为求有多少 $(p,q,z)$ 满足 $l_1p+l_2q+z=K$。\pause

一旦确定了 $p,q$，则唯一确定了 $z$。故答案为求有多少 $(p,q)$ 满足 $l_1p+l_2q\le K$。即为

$$
\sum_{p=0}^t(\left\lfloor\frac{K-l_1p}{l_2} \right\rfloor+1)
$$\pause

复杂度 $O(l_3^2\cdot\log t)$。
\end{frame}

\subsubsection{万能欧几里得算法}

\begin{frame}{万能欧几里得算法}
很多类欧几里得问题都可以归结为以下模型：
\bigskip

在平面直角坐标系内，有一条直线 $y=\dfrac{ax+b}{c}$，作出所有网格线 $x=k(k\in\N^*),y=k(k\in\N^*)$。

在 $(0,n]$ 中从左往右扫描该直线，与一条横的网格线有交点时进行一次 $U$ 操作，与一条竖的网格线有交点时进行一次 $R$ 操作。特别地，若同时与横线竖线有交点（即经过整点），规定先 $U$ 后 $R$。

其中 $U$ 和 $R$ 都是半群元素，即满足结合律。
\end{frame}

\begin{frame}
用 $f(n,a,b,c,U,R)$ 表示当前问题的答案。直线上下平移整数不影响答案，这意味着 $b$ 可对 $c$ 取模。以下默认 $b<c$。\pause

若 $a\ge c$，则每跨过一条竖线，会至少跨过 $\lfloor\dfrac{a}{c}\rfloor$ 条横线。\pause

$y=\dfrac{ax+b}{c}=\dfrac{(a\bmod c)x+b}{c}+\lfloor\dfrac{a}{c}\rfloor x$，故

$$
f(n,a,b,c,U,R)=f(n,a\bmod c,b,c,U,U^{\lfloor\frac{a}{c}\rfloor}R)
$$
\end{frame}

\begin{frame}
若 $a<c$，考虑将直线沿 $y=x$ 对称。记 $m=\lfloor\dfrac{an+b}{c}\rfloor$，一共会产生 $m$ 个 $U$。第 $j$ 个 $U$ 发生在跨过 $y=j$ 时，此时对应的 $x$ 坐标为 $\dfrac{jc-b}{a}$，故第 $j$ 个 $U$ 之前 $R$ 的数量为 $k_j=\lfloor\dfrac{jc-b-1}{a} \rfloor$。\pause

现在整个操作序列被 $m$ 个 $U$ 分割，第一个 $U$ 前有 $k_1$ 个 $R$，最后一个 $U$ 后有 $n-k_m$ 个 $R$，中间可以通过翻转坐标系变成 $a\ge c$ 的版本。故

$$
f(n,a,b,c,U,R)=R^{k_1}\cdot U\cdot f(m-1,c,c-b-1,a,R,U)\cdot R^{n-k_m}
$$
\end{frame}

\subsubsection{}

\begin{frame}{P5170【模板】类欧几里得算法}
给定 $n,a,b,c$，分别求 $\sum\limits_{i=0}^n\lfloor\dfrac{ai+b}{c}\rfloor,\sum\limits_{i=0}^n\lfloor\dfrac{ai+b}{c}\rfloor^2,\sum\limits_{i=0}^n i\lfloor\dfrac{ai+b}{c}\rfloor$。$t$ 组数据。

$t\le 10^5$，$0\le n,a,b,c\le 10^9$，$c\neq 0$。
\end{frame}

\begin{frame}
以 $\sum\limits_{i=0}^n\lfloor\dfrac{ai+b}{c}\rfloor$ 为例。

设计二元组 $(y,ans)$，$U$ 操作会让 $y\leftarrow y+1$，$R$ 操作会让 $ans\leftarrow ans+y$。用矩阵描述 $U$ 和 $R$，显然是满足结合律的。故可直接套用万欧。
\end{frame}

\begin{frame}{QOJ8706 解方程}
求方程 $(a+bx)\bmod p\in[0,c_2),x\in[0,c_1)$ 的所有解。

$a,b$ 随机生成，$p\in[9\times10^{17},10^{18}]$ 且是质数，$c_1,c_2\le 100\sqrt{p}$。
\end{frame}

\begin{frame}
可以发现解的数量不多。对 $x$ 的区间进行分治判定，假设有 $k$ 个解，则将问题转化为 $O(k\log c_1)$ 次判定一个 $x$ 的区间中是否存在解。\pause

将 $(a+bx)\bmod p<c$ 改写为 $bx-p\left\lfloor\dfrac{a+bx}{p} \right\rfloor<c-a$。只需记录在一个 $x$ 的区间中 $S=bx-p\left\lfloor\dfrac{a+bx}{p}\right\rfloor$ 的最小值。\pause

使用万欧的模型，维护二元组 $(S,M)$，直线 $y=\dfrac{a+bx}{p}$ 遇到竖线则 $S\leftarrow S+b,M=\min\{M,S\}$，遇到横线则 $S\leftarrow S-p$。$R$ 与 $U$ 操作可以用 $\min,+$ 矩阵表示。
\end{frame}

\subsection{连分数}

\subsubsection{定义}

\begin{frame}{连分数}
形如
$$
x = a_0 + \cfrac{1}{a_1 + \cfrac{1}{a_2 +\cfrac{1}{\cdots+\cfrac{1}{a_n}}}}
$$
的分数被称为连分数，简记为 $[a_0,a_1,\cdots,a_n]$。若其中每个 $a_i$ 均为正整数，则称为简单连分数。下面所说的连分数均指简单连分数。并定义 $x_k=[a_0,a_1,\cdots,a_k]$ 为 $x$ 的 $k$ 阶渐近分数。
\end{frame}

\subsubsection{有理数表示}

\begin{frame}{有理数表示}
\begin{theo}
一个有理数有且仅有两种连分数表示。
\end{theo}\pause

用辗转相除法构造连分数：令 $r_0=x$，$a_k=\lfloor r_k\rfloor$，$r_{k+1}=\dfrac{1}{r_k-a_k}$，当 $r$ 为整数时停止。\pause

这种方式可以构造出一个连分数，且位数是 $\log$ 级别的。\pause

若最后一位 $a_n\neq 1$，则 $[a_0,a_1,\cdots,a_n-1,1]$ 是另一种连分数表示。若 $a_n=1$，则 $[a_0,a_1,\cdots,a_{n-1}+1]$ 是另一种连分数表示。
\end{frame}

\begin{frame}
\begin{proof}
证明任意有理数不存在大于两种不同的连分数表示。\pause

只需证：若规定连分数的最后一位大于 $1$，则连分数表示唯一。\pause

假设 $x=[a_0,a_1,\cdots,a_n]=[b_0,b_1,\cdots,b_m]$，其中 $a_n>1,b_m>1$。由于 $[a_1,\cdots,a_n]>1$，知 $[a_1,\cdots,a_n]^{-1}<1$，则 $a_0=\lfloor x\rfloor$。同理 $b_0=\lfloor x\rfloor$。\pause

归纳即证。
\end{proof}\pause

一般将 $a_n>1$ 的连分数表示称为标准表示。
\end{frame}

\subsubsection{渐近分数}

\begin{frame}{渐近分数}
对于有理数 $x$，其 $k$ 阶渐近分数为 $x_k=\dfrac{p_k}{q_k}=[a_0,\cdots,a_k]$。直接使用定义计算渐近分数并不方便。一般有如下递推公式：
$$
\left\{
\begin{aligned}
&p_0=a_0,p_1=a_0a_1+1\\
&p_k=a_k\cdot p_{k-1}+p_{k-2}\quad (k\ge 2)\\
&q_0=1,q_1=a_1\\
&q_k=a_k\cdot q_{k-1}+q_{k-2}\quad (k\ge 2)
\end{aligned}
\right.
$$
\end{frame}

\begin{frame}
\begin{proof}
当 $a_{0\sim k-1}$ 确定，$p_k,q_k$ 可以被表示为 $a_k$ 的一次式，设 $p_k=A_ka_k+B_k,$，$q_k=C_ka_k+D_k$。则有：

$$
\begin{aligned}
\frac{p_{k+1}}{q_{k+1}}&=[a_0,a_1,\cdots,a_k+\frac{1}{a_{k+1}}]\\ \pause
&=\frac{A_k(a_k+\frac{1}{a_{k+1}})+B_k}{C_k(a_k+\frac{1}{a_{k+1}})+D_k}\\ \pause
&=\frac{(A_ka_k+B_k)a_{k+1}+A_k}{(C_ka_k+D_k)a_{k+1}+C_k}
\end{aligned}
$$\pause

则 $B_{k+1}=A_k,A_{k+1}=p_k$。于是 $p_k=A_ka_k+B_k=p_{k-1}a_k+A_{k-1}=p_{k-1}a_k+p_{k-2}$。$q_k$ 同理。
\end{proof}
\end{frame}

\subsubsection{}

\begin{frame}{P10788 [NOI2024] 分数}
定义正分数为分子分母都为正整数的既约分数。定义完美正分数集合 $S$ 为满足以下五条性质的集合：

\begin{itemize}
\item $\dfrac{1}{2}\in S$
\item 对于 $\dfrac{1}{2}<x<2$，$x\not\in S$
\item 对于所有 $x\in S$，$\dfrac{1}{x}\in S$
\item 对于所有 $x\in S$，$x+2\in S$
\item 对于所有 $x\in S$ 且 $x>2$，$x-2\in S$
\end{itemize}

求 $1\le i\le n,1\le j\le m$ 的完美正分数 $\dfrac{i}{j}$ 数量。$2\le n,m\le 3\times10^7$。
\end{frame}

\begin{frame}
把分数用标准连分数表示，取倒数相当于在连分数前面增加或删除一个 $0$，$\pm 2$ 相当于对连分数的第一位 $\pm 2$。于是可知：$x\in S\iff x=[2k_0,2k_1,\cdots,2k_n]$，其中 $k_0\ge 0,k_1,\dots,k_n\ge 1$。\pause

由于有理数的标准连分数表示是唯一的，直接 DFS 搜索所有可能的连分数就能得到答案。具体地，从右到左加入连分数的每一位，实时维护其分数表示 $\dfrac{a}{b}$。这样的复杂度为 $O(ans)$，已经可以获得 $90$ 分。\pause

后续的优化方式是考虑枚举连分数表示数列的最大值位置，然后可以一次统计好多个完美正分数。复杂度可能还是 $O(ans)$，但是常数小很多。
\end{frame}

\end{document}
